\chapter{Tinjauan Pustaka}

\section{Penelitian Terdahulu}
\TemplateTodo{Ringkas penelitian terdahulu yang paling relevan. Jelaskan tujuan, metode, data, hasil, dan keterbatasannya. Akhiri dengan posisi penelitian Anda dibandingkan karya-karya tersebut.}

Tinjauan pustaka bukan sekadar daftar ringkasan artikel. Susun pembahasan secara tematik sehingga pembaca memahami perkembangan pengetahuan, kesenjangan yang masih ada, dan alasan penelitian ini diperlukan.

Ringkasan penelitian pembanding disajikan pada Tabel \ref{tab:penelitian-terdahulu}. Sebutkan tabel di dalam kalimat sebelum tabelnya muncul, karena pembaca perlu tahu apa yang akan dibacanya.

% CONTOH TABEL 1: seluruh kolom lentur, memakai bobot L{...}.
%
% Tabel ringkasan penelitian terdahulu biasanya panjang, sehingga dipakai
% xltabular: lebarnya menyesuaikan seperti tabularx, tetapi isinya boleh
% terpotong antarhalaman seperti longtable. Karena itu tabel TIDAK dibungkus
% environment table dan tidak memakai [H]: xltabular mengatur posisinya sendiri.
%
% Baris judul di antara \endfirsthead dan \endhead diulang pada setiap halaman
% lanjutan, jadi pembaca tidak perlu kembali ke halaman pertama untuk tahu arti
% tiap kolom. \endlastfoot menutup tabel dengan satu \hline, sehingga baris data
% terakhir tidak perlu diberi \hline lagi.
%
% Kolom L{bobot} adalah kolom X berbobot yang rata kiri. Jumlah seluruh bobot
% harus sama dengan cacah kolomnya: lima kolom di bawah ini berjumlah
% 0,70 + 1,05 + 1,20 + 1,15 + 0,90 = 5,00. Kolom yang isinya panjang diberi
% bobot lebih besar. Definisinya ada di config/thesis-preamble.tex.
\begin{xltabular}{\linewidth}{|L{0.70}|L{1.05}|L{1.20}|L{1.15}|L{0.90}|}
\caption{Ringkasan penelitian pembanding langsung.}\label{tab:penelitian-terdahulu}\\
\hline
Penulis (Tahun) & Metode & Dataset & Hasil & Keterbatasan \\
\hline
\endfirsthead
\hline
Penulis (Tahun) & Metode & Dataset & Hasil & Keterbatasan \\
\hline
\endhead
\hline
\endlastfoot
Nama peneliti (tahun) \cite{exampleReference} & Ringkasan pendek metode yang dipakai & Nama dan ukuran dataset beserta cakupannya & Angka hasil utama yang dilaporkan & Celah yang ditanggapi penelitian Anda \\
\hline
Nama peneliti (tahun) & Tambahkan satu baris untuk setiap penelitian yang ditinjau & Sebutkan sumber datanya & Nyatakan hasil dengan satuan yang jelas & Tuliskan keterbatasan secara spesifik \\
\end{xltabular}

\section{Landasan Teori}
\TemplateTodo{Jelaskan konsep, teori, algoritma, standar, atau kerangka kerja yang benar-benar digunakan pada penelitian.}

\subsection{Konsep Utama Pertama}
Tuliskan definisi, prinsip kerja, dan relevansi konsep terhadap penelitian. Gunakan sitasi untuk definisi atau teori yang bukan pendapat pribadi.

\subsection{Konsep Utama Kedua}
Tambahkan konsep lain sesuai kebutuhan. Jika menggunakan rumus, jelaskan arti setiap simbol dan alasan rumus tersebut digunakan.

\section{Kerangka Pemikiran}
\TemplateTodo{Jelaskan hubungan antara masalah, teori, metode, data, dan keluaran penelitian. Bagian ini dapat dibantu dengan gambar alur konseptual.}

\begin{figure}[H]
\centering
\includegraphics[width=0.75\linewidth]{example-figure.pdf}
\caption{Contoh penggunaan gambar pada template}
\label{fig:example-figure}
\end{figure}

Contoh penggunaan gambar ditunjukkan pada Gambar~\ref{fig:example-figure}. Notasi matematika dapat digunakan untuk menjelaskan hubungan antar lapisan pada Gambar~\ref{fig:example-figure}. Misalkan vektor masukan dinyatakan sebagai
\begin{equation}
\mathbf{x} =
\begin{bmatrix}
x_1 & x_2 & \cdots & x_n
\end{bmatrix}^{\mathsf{T}}.
\end{equation}
Untuk lapisan tersembunyi ke-$l$, aktivasi neuron dihitung dengan
\begin{equation}
\mathbf{a}^{[l]} = \sigma\left(\mathbf{W}^{[l]}\mathbf{a}^{[l-1]} + \mathbf{b}^{[l]}\right),
\qquad l = 1,2,\ldots,L-1,
\end{equation}
dengan $\mathbf{a}^{[0]}=\mathbf{x}$, $\mathbf{W}^{[l]}$
adalah matriks bobot pada lapisan ke-$l$,
$\mathbf{b}^{[l]}$ adalah vektor bias, dan
$\sigma(\cdot)$ adalah fungsi aktivasi. Keluaran model dinyatakan sebagai
\begin{equation}
\hat{\mathbf{y}} = f\left(\mathbf{W}^{[L]}\mathbf{a}^{[L-1]} + \mathbf{b}^{[L]}\right),
\end{equation}
dengan komponen keluaran
\begin{equation}
\hat{\mathbf{y}} =
\begin{bmatrix}
\hat{y}_1 & \hat{y}_2 & \cdots & \hat{y}_k
\end{bmatrix}^{\mathsf{T}}.
\end{equation}

\subsection{Gambar dengan Keterangan Panel}
Gambar yang terdiri atas beberapa panel diberi keterangan (a), (b), dan seterusnya. Alih-alih membuat beberapa environment \texttt{figure} yang masing-masing mendapat nomor sendiri, satu gambar dibungkus \texttt{tikzpicture} lalu keterangan panel ditulis sebagai simpul di bawahnya, sehingga seluruh panel tetap berbagi satu nomor dan satu caption.

\begin{figure}[H]
\centering
% Sintaks ($A!0.25!B$) memerlukan pustaka calc dan berarti "titik pada 25%
% jarak dari A ke B", jadi keterangan panel selalu berada tepat di bawah
% pusat panelnya berapa pun lebar gambarnya.
\begin{tikzpicture}
  \node[anchor=south west, inner sep=0pt] (img) at (0,0)
    {\includegraphics[width=0.85\linewidth]{example-figure.pdf}};
  \node[anchor=north, font=\small, yshift=-2mm]
    at ($(img.south west)!0.25!(img.south east)$) {(a) sebelum pemrosesan};
  \node[anchor=north, font=\small, yshift=-2mm]
    at ($(img.south west)!0.75!(img.south east)$) {(b) sesudah pemrosesan};
\end{tikzpicture}
\caption{Contoh gambar dengan dua keterangan panel}
\label{fig:contoh-panel}
\end{figure}

Perbandingan kedua kondisi ditunjukkan pada Gambar~\ref{fig:contoh-panel}. Untuk gambar yang seluruhnya digambar dengan TikZ atau pgfplots, kerjakan di folder \texttt{tikz/} lalu sisipkan berkas PDF hasilnya; lihat \texttt{tikz/README.md}.
